\section{Design}
\label{sec:design}

\houdini is built using Python 3 and leverages a QEMU VM to provide a controlled testing environment for container security validation. We use a QEMU-based VM for two reasons: accurate execution environment and cross-architecture support. First, QEMU supports hardware-assisted virtualization, providing a near-native execution environment that closely mirrors real-world behavior. Second, when hardware support is unavailable, QEMU enables full system emulation, allowing tests to be conducted on architectures different from the host CPU.

Instead of running containers directly on the host system, \houdini starts a dedicated VM, ensuring that any security tests remain isolated and do not affect the host system. Dedicated VMs also help ensure test reproducibility. Within the \houdini configuration, specific versions of Docker, the Linux kernel, and runc can be requested and installed into the VM\footnote{While \houdini currently only works with Docker containers, we expect that other container runtimes would be straightforward to support given that most have some degree of Docker compatibility.}. The VM is created with the help of Buildroot\cite{buildroot}, which enables the construction of minimal, reproducible Linux environments. 
%tailored for security testing. 
This granularity in configuration enables our goal of easily testing how different kernel and container runtime versions impact policy enforcement. It also enables regression testing across versions to detect when security behavior changes unexpectedly.

At the core of \houdini's architecture (see Figure~\ref{fig:state-machine}), a Python Flask-based server runs on the host operating system. This server acts as a bridge between the host and the VM, facilitating communication between them over a lightweight HTTP API. The Flask server is responsible for sending commands and configuration data to the VM and receiving trick execution results. When a trick is initiated, the \houdini client on the host instructs the VM to start a new trick. The VM then launches a container with the specified security configurations and executes the trick inside it.

When a trick is initiated, the Flask server exposes an endpoint that the \houdini client on the host uses to send a request containing the container configuration and the test logic to execute. The server then transmits this data into the VM. Inside the VM, a listener service receives this request, parses the configuration, and launches a container with the specified settings. The trick code is then executed within the container. Once execution is complete, the listener inside the VM collects the results, and sends them back to the Flask server. The Flask server then responds to the original request with the results, allowing the host to log the outcome.

If \houdini were compromised during testing, it would not necessarily pose a threat to the host system due to the design of the testing environment. \houdini is intentionally designed to operate within a controlled, containerized environment inside a VM. This means that even if a trick were to exploit a vulnerability within the containerized environment, the impact would be contained within the VM, isolated from the host system and other containers. While the VM's primary purpose here is to provide a consistent testing environment, it also serves as an additional layer of security.

\begin{figure}[h]
  \includegraphics[width=1\linewidth]{figs/designNew.pdf}
  \caption{Architecture of \houdini}\label{fig:state-machine}
\end{figure}
